\chapter{Codifica}
\label{cap:codifica}
Dopo aver completato la fase di progettazione è stata eseguita la fase di codifica e validazione. Questo capitolo sarà il conseguimento della seconda parte del periodo di stage dove verranno spiegate le funzionalità implementate nella Skill mostrando porzioni di codice. Verrà inoltro spiegata la struttura del progetto: le varie cartelle contenti i file delle \textit{classi} realizzate e i \textit{packages} installati. 

\section{Organizzazione del codice}
Successivamente alla parte di comprensione e configurazione dei servizi utilizzati in questo progetto, e alla creazione della Skill, è stato eseguito un micro-processo di organizzazione del codice prima della sua stesura. Questa pratica adottata è stata necessaria per avere uno schema e una gerarchia delle parti di codice che sono state sviluppate durante il periodo di tirocinio. Tale organizzazione ha permesso di sviluppare in maniera più ordinata e consapevole riducendo la probabilità di introdurre errori e rendendo il codice più manutenibile. Pertanto la cartella del progetto è stata così disposta: